\begin{abstract}
\noindent
Static-prefix multi-agent LLM coding systems route work across a
cheap-to-expensive model-tier hierarchy (haiku / sonnet / opus) to
manage cost. Existing routing research targets the
\emph{subagent-dispatch} problem and produces clear (role
\(\times\) signal) rules; the analogous \emph{main-thread
orchestrator} problem has no role taxonomy and requires per-turn
classification, but it dominates the cost shape. We measure
\blackrim's production cost split as 99.4\,\% main-thread / 0.6\,\%
subagent dispatch, indicating that the orchestrator itself is the
binding cost lever in a well-tuned multi-agent system. We present
two complementary primitives that target this lever: (i) a
semantic-similarity per-turn router for the main thread, drawing
on FrugalGPT cascading and RouteLLM-style preference-data routing;
(ii) an agentic plan-level semantic cache stacking on top of
existing prefix caching, drawing on the 2025 NeurIPS
\emph{Agentic Plan Caching} result. The router is grounded in
exemplar curation across \(\sim\)150 hand-labeled turns (50 per
tier) and a k-NN classifier with conservative escalation to opus
when similarity falls below threshold; the plan cache is keyed on a
semantic signature of the orchestration plan and indexes
strategy-level reuse across structurally-similar bd dispatches.
We report a 4.56\(\times\) read-to-creation ratio on the existing
prefix cache (69.3\,\% baseline savings) as the calibration point
the new techniques stack atop. We also document an integration
blocker: \texttt{Claude Code} currently exposes only session-level
model selection, not per-turn routing; we evaluate three workaround
paths (Agent SDK middleware, Claude Code feature request, prompt
injection). Closing the loop, we discuss why \blackrim's prior
subagent-dispatch advisor (CC-TS, \cite{moa1landscape}) does
\emph{not} transfer directly to the main thread and what would.
\end{abstract}
